\section{Related Work}
\label{sec:related}

\para{Long and irrelevant context.}
A line of work shows that language models get worse as their input grows.
\citet{shi2023distracted} add irrelevant sentences to math word problems and find large accuracy drops.
\citet{liu2024lost} show that models use information at the start and end of a long context better than information in the middle.
\citet{levy2024same} pad reasoning inputs with filler and find that accuracy falls well before the context limit.
These studies vary the \emph{input} that the model reasons over.
Unlike them, we vary the \emph{task instruction} and hold the input fixed, and we work at realistic prompt lengths (about 100 to 1{,}700 tokens) rather than multi-thousand-token contexts.
Our padding control borrows the logic of filler-based studies, but we use it to separate a length effect from a content effect inside the instruction.

\para{Instruction following under many constraints.}
IFEval~\citep{zhou2023ifeval} measures whether models obey verifiable format constraints, and FollowBench~\citep{jiang2024followbench} shows that compliance drops as constraints are added one level at a time.
Our nested design is similar in spirit, since each tier adds content to the previous one.
Unlike these benchmarks, our tiers elaborate a \emph{single} fixed format constraint rather than add new ones, so we can ask whether extra explanation helps or hurts compliance with the same requirement.
We also measure task correctness separately from compliance.

\para{Prompt sensitivity and instruction wording.}
Models are sensitive to surface features of prompts, including formatting~\citep{sclar2024quantifying} and whether the instruction is even meaningful~\citep{webson2022prompt}.
\citet{mishra2022reframing} find that reframing long, complex instructions into shorter, itemized ones improves GPT-3 performance.
\citet{zhao2021calibrate} show that few-shot predictions are biased toward labels that are frequent or appear near the end of the prompt.
The label collapse we observe in the 0.5B model resembles this recency and majority-label bias, but it is caused by instruction rules rather than demonstrations.
Prompt compression methods such as LLMLingua~\citep{jiang2023llmlingua} assume that long prompts can be shortened with little loss; our results suggest that, for capable models, shortening also brings little gain.

\para{Instruction-tuned models.}
Instruction tuning~\citep{wei2022finetuned,ouyang2022training} trains models to follow natural-language task descriptions.
We study one modern instruction-tuned family, \qwen~\citep{qwen2025qwen25}, across a within-family scale ladder, which lets us attribute differences to model size rather than training recipe.
